\section{Introduction}
\label{sec-1}

\begin{comment}
\textit{
\begin{itemize}
    \item STATEMENT: Initial paragraph stating the problem that the paper addresses and its proposed solution including its novelty and its applications. {\bf DONE}
    \item SOTA: Give a concise description of the state of the art of Chebyshev accelerated subspace iteration with many references to scientific papers, HPC optimizations, applications to domain software and stand alone libraries. {\bf DONE}
    \item SUBSECTION: brief description of the general subspace iteration + Chebyshev filtering and RR-projection. Introduce the ChASE library specialized algorithm, its optimizations (polynomial degree, spectral bounds), its emphasis on HPC and parallelism. {\bf DONE} 
    \item SUBSECTION: Focus on QR-factorization. Enlist the many algorithms for factorizing tall and skinny matrices. Describe the stability and the performance aspects of each of them. Emphasize the issues that accompany the CholeskiQR and how they have been addressed in the last 10 years. Describe how the Chebyshev filtering may make the use of CholeskiQR unfeasible if not appropriately controlled. State in details the goal of the paper and the results that achieves. State how such results can be used to switch from different QR-algorithms including the various flavors of CholeskiQR and how the switching mechanism can be used to enhance performance in ChASE. {\bf DONE}
    \item Describe the structure of the manuscript and its original contributions. {\bf DONE}
\end{itemize}
}
\end{comment}

Subspace iteration with Chebyshev acceleration is among the earliest and most influential iterative techniques developed for computing a selected portion of the spectrum of symmetric (Hermitian) algebraic eigenvalue problems, and it remains a cornerstone of modern large-scale eigensolvers. A trial subspace is first enhanced through a polynomial filtering step that amplifies the components associated with the desired eigenvalues, after which the resulting vectors are orthonormalized and used to construct a reduced eigenvalue problem via a Rayleigh–Ritz projection. This reduced problem is then solved using standard dense diagonalization techniques (e.g., divide-and-conquer algorithm), and the resulting approximate eigenpairs are projected back to the original space and tested for convergence. 

From a computational standpoint, this workflow is dominated by two key computations: the filtering step, which can be efficiently expressed as repeated matrix–matrix multiplications (GEMM) and therefore maps well to modern hardware, and the orthonormalization step, which is typically performed using the Householder QR algorithm and constitutes the second most expensive operation. Unlike the filter, Householder QR relies heavily on lower-level BLAS routines and is less amenable to efficient high-performance implementations. A natural alternative is the CholeskyQR algorithm, which is rich in BLAS-3 operations and can be implemented in a manner similar to the Chebyshev filter. However, CholeskyQR is notoriously sensitive to the condition number of the matrix of vectors to be orthonormalized, a difficulty exacerbated by the fact that polynomial filtering tends to produce highly correlated subspaces. Recent variants such as CholeskyQR2 \cite{fukaya2014choleskyqr2} and shifted 
CholeskyQR2 \cite{fukaya2020shifted} offer promising alternatives, provided reliable information about the condition number of the subspace is made available. In this work, we present a detailed numerical analysis of the conditioning of the filtered subspace under an exhaustive range of scenarios and show how this information can be exploited to safely modify the 
orthonormalization step within subspace iteration. We implement the resulting strategy in 
the ChASE eigensolver library and demonstrate that it delivers both numerical robustness and significant performance gains on modern architectures.


\subsection{Filtered subspace iteration: its inception and modern revival}
One of the the first explicit mention in the scientific literature of
subspace iteration applied to the solution of the symmetric algebraic
eigenvalue problems is by L.~Bauer in
1957~\cite{Bauer:1957hv}. Following this work, as famously mentioned
by Parlett~\cite{Parlett:1998tv}, eigensolvers based on subspace
iteration [ ... where developed in the 1960s and 1970s when the
Lanczos algorithm was under the cloud]. What shaped subspace
iteration in its modern form is the fundamental work of Rutishauser in
a number of papers spanning from 1969 to
1971~\cite{Rutishauser:1969ub, Rutishauser:1971wc,
  Rutishauser:1970wc}. The first paper of this
series~\cite{Rutishauser:1969ub} builds on Bauer's Simultaneous
Iteration method and introduce many of the key ideas included in
modern implementations of eigensolvers based on subspace
iteration. Soon after is groundbreaking work, Rutishauser published
two similar papers with additional details of the algorithm (named {\it
  ritzit}) which are at the base of its robustness and
efficacy~\cite{Rutishauser:1970wc, Rutishauser:1971wc}.

Soon after Rutishauser's first publication, several authors
contributed to improve his initial algorithm. For instance, the Ritz
iteration, which projects the eigenproblem onto the search space using
a Jacobi step, was upgraded by Stewart to a Rayleigh-Ritz step so as
to avoid to deal with non-positive definite
matrices~\cite{Stewart:1969cq}. Stewart rigorously
demonstrated the enhanced convergence of the eigenpairs when
orthogonal iteration (subspace + orthonormalization) is used in
conjunction with Rayleigh-Ritz. Concurrently to these development a
version of subspace iteration was developed by civil engineers
proposing an algorithm very similar to Rutishauser's, but their
exposition is less clear-cut and uses a language and formalism
somewhat different from conventional numerical
analysts~\cite{Clint:1970wz}. The interested reader can find a clear
review of the improved {\it ritzit} algorithm in Parlett
book~\cite{Parlett:1998tv}.

Alongside these developments, several authors generalized subspace
iteration to non-Hermitian and unsymmetric eigenproblems
(see~\cite{Stewart:1976vaa} for an example). In the intervening years
up to the last decade, the development of iterative eigensolvers for
the Hermitian eigenvalue problem took a different direction due to the
resurgence of the Lanczos algorithm and its
variations~\cite{Cullum:1974fe, Parlett:1979vf, Cullum:1985wj,
  Wu:2000vq, Davidson:1975wf, Crouzeix:1994wh, Ruhe:1984ep,
  Stewart:2002tw}. Subspace iteration eigensolvers saw a revival in
popularity starting in the middle of the 2000s with application to
electronic structure theory, first in the context of sparse
eigenproblems~\cite{Zhou:2006ek, Zhou:2014fe}, and later also for
sequences of dense eigenproblems~\cite{Berljafa:2014jv, Levitt:2015wc}.
This coming back is due in part to the emerging need of solving for
just the full subspace without resolving the single
eigenpairs~\cite{Zhou:2006ek} and in part to the ability of the
eigensolver to receive as input approximate eigenvectors, and in doing
so, decrease considerably its time-to-convergence~\cite{DiNapoli:2012uy}.

The last decade has witnessed an increased use of filtered subspace iteration in new advanced methods of electronic structure computations. A Chebyshev filtered subspace iteration has been used to power a discontinuous Galerkin approach to solve the Kohn-Sham equation using a adaptive local basis set \cite{Banerjee2016, Banerjee2018}. Similarly, an early version of the ChASE library \cite{Winkelmann2019chase} was integrated to solve Excitonic Hamiltonian emerging from the Tamm-Dancoff approximation (TDA) of the Bethe-Salpeter equation \cite{zhang2021solving}. It has been also shown that Chebyshev Filtered subspace returns high-performance gains on many cores \cite{Kreutzer2018}. More recently, our team has endevoured in a substantial effort to develop the ChASE library into a full-fledged accelerated and distributed parallel library \cite{Wu2022chase,Wu2023chase} with further extensions to pseudo-Hermitian eigenproblems \cite{dinapoli2026chase} arising from Bethe-Salpeter equation beyond TDA.      


%\subsection{Quotient iteration and  Raleigh-Ritz (R-R) projection.}
%\label{sec-2-1}
% {\it In this section we need to mention convergence theorems
%   only. There is no need mention backward stability. Sensitivity of
%   the solution to starting vectors can be just mentioned as an element
%   of stability already suggested by Rutishauser. No theorems needed
%   for that.
% }


\begin{comment}
For sake of simplicity, and only in this section, we make the two
following additional hypothesis: first we assume that $A$ is positive
definite ($\lambda_1 > 0$), second we seek the highest $k$ eigenpairs
corresponding to $(\lambda_{n-k},x_{n-k}),\ldots,(\lambda_n,x_n)$.
Algorithm~\ref{alg:SubIt} illustrate the traditional structure of the
Quotient Iteration algorithm including a Rayleigh-Ritz step. This is
by no means the most general subspace iteration algorithm but it is
complex enough to aid in analyzing its conventional convergence
properties.

\newlength{\dlen} \settowidth{\dlen}{{\sc Check for convergence}} 
\newlength{\ylen} \settowidth{\ylen}{Y}
\begin{algorithm}[h!t]
  \caption{Quotient Iteration + Rayleigh-Ritz}
  \label{alg:SubIt}
  \begin{algorithmic}[1]
  \Require Matrix $A$, % vector matrix 
  % $ \subit{V}{0} \equiv \left[
  %   v^{(0)}_1, \ldots, v^{(0)}_{k} \right]$
  \Ensure $k$ extremal eigenpairs $\left(\Lambda,Y\right)$ with
  $\Lambda = \diag (\lambda_n, \ldots,\lambda_{n-k})$ and $Y = \left[
    y_n, \ldots, y_{n-k}\right]$.
  \item[]
  \State Set the size of search space $\ell > k$ 
  \State $\subit V{0} \gets {\rm randn}(n,\ell)$
  \While{size$(Y) < k$}
  \State Compute, $\subit V{i}= A \subit V{i-1}.$ \label{lst:line:subsp} 
  \Comment{\parbox{\dlen}{\sc Subspace Projection}}
  \State Orthonormalize $\subit V{i} = \subit Q{i} \subit{R}{i}.$ \label{lst:line:QR}
  \Comment{\parbox{\dlen}{\sc QR algorithm}} 
  \State Compute Rayleigh quotient $\tilde{A} = \her{\subit Q{i}} A\
  \subit Q{i}.$ \label{lst:line:rrst} 
  \Comment{\parbox{\dlen}{{\sc Rayleigh-Ritz} (Start)}}
  \State Solve the reduced problem $\tilde{A} =\subit W{i}
  \subit{\tilde{\Lambda}}{i} \her{\subit W{i}}$.
  \State Compute $\subit V{i} \gets \subit Q{i} \subit W{i}.$ \label{lst:line:rren}
  \Comment{\parbox{\dlen}{{\sc Rayleigh-Ritz} (End)}}
    \If{${\rm Res}(\subit{V_{:,a}}{i},\tilde{\lambda}_a)<\textsc{tol}$}
    \label{lst:line:resid} 
    \Comment{\parbox{\dlen}{{\sc Check for convergence}}}
    \State $\parbox{\ylen}{$\Lambda$}= \big[\Lambda\;\tilde{\lambda}_a\big]$
    and $\parbox{\ylen}{Y}= \big[Y\;\subit{V_{:,a}}{i}\big]$ 
    \EndIf
  % \State Compute residuals ${\rm Res}(\subit{V_{:,j}}{i},\tilde{\lambda}_j)<\textsc{tol}$
  \EndWhile
\end{algorithmic}
\end{algorithm}

The simplest version of subspace iteration is realized by removing
from Algorithm~\ref{alg:SubIt} all the lines between
Line~\ref{lst:line:subsp} and~\ref{lst:line:resid}. Notice that we are
already considering a search space which is larger than the sought
after subspace. Such a choice increases the convergence ratio at the
cost of extra computations. The resulting subspace converges linearly
to the desired eigenspace under the mild assumption that the matrix
$\her{X} \subit V{0}$ is indeed an $\ell\times \ell$ invertible
matrix. Let us define with
$\subit \Scal{i} = \Rcal\left(\subit V{i}\right)$ the subspace spanned
by the column of $\subit V{i}$ and with $\Theta^{(i)}$ the largest
canonical angle between $\subit \Scal{i}$ and $\Rcal (X_1)$ with
$X_1 = [x_n, \ldots , x_{n-\ell}]$. Then $\exists\ \varepsilon_i > 0$
such that
% In particular, for any
% $\varepsilon > 0$ one can define $\Theta^{(i)}$ as the largest canonical
% angle between $\Rcal (A \subit V{i-1})$ and $\Rcal (X_1)$ with
% $X_1 = [x_n, \ldots , x_{n-k}]$, such that
\[
\Theta^{(i)} = \odg{\left[ \frac{\lambda_{n-\ell-1}}{\lambda_{n-k}} +
    \varepsilon_i \right]^i}
\]
holds (see~\cite{Stewart:2001id} pp.389-90) with $\lim_{i
  \longrightarrow \infty} \varepsilon_i = 0$. This relation
illustrates the convergence of whole subspace and as such it does not give any
indication on how well each single eigenvector $x_a$ can be
approximated by $\subit \Scal{i}$. A refined version of the
convergence bound for a single vector is given by
\[
\tan \left( \theta_a^{(i)}\right) \leq \left[ \frac{\lambda_{n-\ell-1}}{\lambda_{a}}
\right]^i \tan \left( \Theta^{(0)}\right) 
\]
where $\theta_a^{(i)}$ is the angle between $x_a$ and
$\subit \Scal{i}$ (see~\cite{Parlett:1998tv} pp.332-3).

When the QR factorization is added to the picture, by adding
Line~\ref{lst:line:QR} to Algorithm~\ref{alg:SubIt}, one can formulate
the convergence to single eigenvectors in a more compact and effective
manner 
\begin{theorem}
\label{th:convsaad}
Let $\subit \Pcal{i}$ and $P_1$ being the orthogonal projectors onto the
subspaces $\subit \Scal{i}$ and $\Rcal (X_1)$ respectively and assume
that $\her{X} \subit V{0}$ is invertible. Then for each eigenvector
$x_a$ of $A$ with $a=1,\dots ,\ell$ there exists a unique vector $v_a$
in $\subit \Scal{0}$ such that $P_1 v_a = x_a$ and the following
inequality is satisfied
  \[
  \left\|\left(I - \subit \Pcal{i} \right)x_a\right\| \leq \|x_a - v_a
  \| \left( \left|\frac{\lambda_{n-\ell-1}}{\lambda_{n-k}}\right| +
    \varepsilon_i \right)^i
  \]
\end{theorem}
(see \cite{Saad:2011tu} pp.119-20 for a demonstration). This theorem
guarantees the existence and uniqueness of the solution within the
subspace $\subit \Scal{i}$ and at the same time bounds from above the
ratio of convergence (in the case where $A$ is diagonalizable a bound
for sequence of $\varepsilon_i$ can also be found). In practice the
computation of the actual subspace is benefitted by the QR
factorization which eliminates rank deficient occurrencies.  It is
worth mentioning that the orthonormalization is already introduced by
Rutishauser~\cite{Rutishauser:1969ub}. In particular, he suggests to
combine few iterations together before performing the
orthonormalization, so as to avoid eccessive computation and extract
the most benefit out of it. As already mentioned, the same paper
introduces already a Ritz iteration which is performed on the product
$\her{\subit V{i}} \subit V{i} = \her{\subit{R}{i}} \subit{R}{i}$.
While this trick reduces the overall amount of floating point
operations, it has the limiting feature of being possible only for
positive definite $A$.

A more general, even if more expensive, alternative is to perform a
Raleigh-Ritz (R-R) step as introduced in~\cite{Stewart:1969cq}. In
this paper the author shows how such a step enhances the
convergence of the quotient iteration in the case of Hermitian
eigenproblems. Through backward perturbation analysis one can infer
that the Ritz vector $w_a$ maintains a linear convergence
\[
\sin \angle \left(x_a, \subit Q{i} \subit w{i}_a\right) = \odg{\theta_a^{(i)}}
\]
while the Ritz value $\tilde{\lambda}_a$ has the additional feature of converging
quadratically 
\[
\left| \subit{\tilde{\lambda}}{i}_a -\lambda_a \right| = \odg{\left[\theta_a^{(i)}\right]^2}
\]
as long as the uniform separation condition is satisfied
(\cite{Stewart:2001id} pp. 286-90). In addition it can be shown that
R-R minimizes the residual of a computed eigenpair and that if such a
residual is small then the computed eigenpair is shown to be an exact
eigenpair of the perturbed matrix $(A + E)$ (\cite{Stewart:2001id}
p.61-3). For the latter reasons the standard stopping criterion for a
subspace iteration is based on the computation of eigenpair residuals
(see Line~\ref{lst:line:resid}) and not on the error. Historically
Rutishauser uses a termination criterion based on the stagnation of
the angles $\theta_a^{(i)}$ and does not compute the residual before
the angle is accepted. While this criterion ensures that the algorithm
will terminate even when the requested accuracy cannot be achieved, it
limits the its flexibility when lower accuracies are required.
{\bf need some refining here. Read Parlett book on R-R and perhaps add
or modified material.} 

\end{comment}
% For
% instance, introduces a re-orthogonalization of the vectors spanning
% the approximate subspace, shows how projecting the eigenproblem onto
% this subspace accelerate convergence of eigenvectors and use a filter
% based on Chebyshev polynomials to further accelerate convergence.
% In this publication some additional technical
% details are introduced such as a search space larger than the required
% subspace, and strategy to sombine iteration so as to avoid eccessive
% re-orthogonalization.

 
% \begin{remunerate}
% \item Convergence theorems for subspace iteration (Stew p.57 [1
%   vector], pp.389-90 [subspace]), (GvL pp.410-11), (Rutishauer 1969
%   although the demonstrations here are rather simplistic and unclear),
%   (Saad theorem 5.2). Stewart makes a spectral decomposition
%   subdividing the spetrum in 3 parts: the wanted spectrum, the extra
%   search space and the unwanted spectrum. The demonstration gives just
%   the order of magnitude of the angle between computed subspace and
%   eigenspace without QR re-orthogonalization.
% \item Golub-vanLoan gives a more complex demonstration including QR
%   but using only a spectral decomposition in two parts. Rutishauser
%   also introduces QR by suggesting that it needs to be done only every
%   so often. Refer to Stewart paper (1969)
% \item Rayleigh-Ritz enhances convergence for eigenvalues and
%   eigenvectors for Hermitian eigenproblems. An erlier verion of Ritz
%   projection was also proposed by Rutishauser (1969, 1970). One can
%   refer to Stewart paper on Non-Hermitian matrices of (1976) which
%   introduces a slightly more sophisticated R-R using Schur
%   complement. Look also at Parlett book (which described Ritzit
%   anyway).If \( \sin(\theta) \) is the angle between X and the
%   computed eigenvectors \( \tilde{X} \) then (Stew pp.286-90):
% \begin{description}
% \item[Eigenvalues] \( \tilde{\Lambda} \) converge with a rate
%                          proportional to \( \mathcal{O}(\theta^2) \)
% \item[Eigenvectors] \( \tilde{X} \) converge with a rate proportional
%                           to \( \mathcal{O}(\theta) \)
% \end{description}
% \item Convergence criterion makes use of residuals (instead of
%   solution error as used by Rutishauser 1969 and 1970).
% \begin{description}
% \item[The Rayleigh quotient] minimize the residual of a computed vector (Stew p.63)
% \item[Backward analysis] if residual is small then the computed
%              eigenvector is shown to be an exact eigenvector of (A +
%              E) (Stew p.61)
% \end{description}
% \item Role of starting vectors (mentioned already in Rutishauser
%   1969). This can be related to the convergence effectiveness and
%   later on to motivate the importance of approximate vectors inputed
%   to the algorithm (Stew p.60)
% \end{remunerate}





\subsection{A modern eigensolver based on Chebyshev Accelerated
  Subspace iteration} 
\label{sec-1-4}
Let $A$ be a symmetric or complex Hermitian matrix of size $n$, having $n$
real eigenvalues ordered such that the relation
$\lambda_1\leq \lambda_2\leq \ldots \leq\lambda_n$ holds. The
ascending ordering, which is the opposite of what is traditionally
found in the literuture, is inspired by application to electronic
structure calculations but it is otherwise equivalent to the more
common descending ordering description. We are interested in the first 
{\sf nev} eigenpairs---eigenvalues and their corresponding eigenvectors
$(\lambda_1,x_1),(\lambda_2,x_2),\ldots,(\lambda_\nev,x_\nev)$---for $\nev\ll n$.
The choice for the wanted part of the eigenspectrum is arbitrary but
is in line with the examples presented in this work. Let us just
remark that subspace iteration with a polynomial filter can be used
equivalently for the highest part of the spectrum and not only to
compute the smallest eigenpairs.

\newlength{\dlen} \settowidth{\dlen}{{\sc Check for convergence}} 
\newlength{\ylen} \settowidth{\ylen}{$Y$}
\newlength{\clen}
\settowidth{\clen}{\small{\sc Deflation \& Locking} (Start)}

\begin{algorithm}[h!t]
  \caption{ChASE: Chebyshev Accelerated Subspace iteration Eigensolver}
  \label{alg:ChASE}
  \begin{algorithmic}[1]
  \Require Matrix $A$, {\it (optional) vector matrix 
  $ \subit{V}{0} \equiv \left[
    v^{(0)}_1, \ldots, v^{(0)}_{\nev} \right]$} 
  \Ensure $\nev$ extremal eigenpairs $\left(\Lambda,Y\right)$ with
  $\Lambda = \diag(\lambda_1, \ldots,\lambda_\nev)$ and $Y = \left[
    y_1, \ldots, y_{\nev}\right]$.
  %\item[]
  \State Set the size of search space $\ell > \nev$ 
  \State Initialize $\subit V{0} \gets {\rm randn}(n,\ell)$ \label{lst:line:init} %\left[ \subit{V}{0}, {\rm randn}(n,\ell-k) \right]$
  \State Estimate $\lambda_1$, $\lambda_{\ell}$ and $\lambda_n$. \label{lst:line:lanczos} 
  \Comment{\parbox{\clen}{\sc Lanczos}}
  \While{size$(Y) < \nev$}
  \State Filter the vectors, $\subit V{i} \gets \p i (A) \subit V{i-1}.$ \label{lst:line:cheby} 
  \Comment{\parbox{\clen}{\sc Chebyshev filter}}
  \State Orthonormalize $\subit Q{i} \gets \big[Y\; \subit V{i}\big] \inv{\subit{R}{i}}.$ \label{lst:line:QR} 
  \Comment{\parbox{\clen}{\sc QR algorithm}} 
  \State Compute Rayleigh quotient $\tilde{A} = \her{\subit Q{i}} A
  \subit Q{i}.$ \label{lst:line:rrstarts} 
  \Comment{\parbox{\clen}{{\sc Rayleigh-Ritz} (Start)}}
  \State Solve the reduced problem $\tilde{A} \subit W{i} =\subit W{i}
  \subit{\tilde{\Lambda}}{i}$.
  \State Compute $\subit V{i} \gets \subit Q{i} \subit W{i}.$ \label{lst:line:rrends}
  \Comment{\parbox{\clen}{{\sc Rayleigh-Ritz} (End)}}
  \For{$a=\textsf{converged} \to \nev$ }  
    \If{${\rm Res}(\subit{V_{:,a}}{i},\tilde{\lambda}_a)<\textsc{tol}$}
    \Comment{\parbox{\clen}{{\sc Residuals}}}
    \State $\parbox{\ylen}{$\Lambda$} \gets \big[\Lambda\;\tilde{\lambda}_a\big]$
    \State $\parbox{\ylen}{$Y$} \gets \big[Y\;\subit{V_{:,a}}{i}\big]$ 
    \Comment{\parbox{\clen}{{\sc Deflation \& Locking}}} \label{lst:line:dlstarts}
    \EndIf
    \EndFor 
    \State $\subit V{i} \gets \big[\subit{V_{:,{\sf nonconverged}: \ell}}{i}\big]$ \label{lst:line:dlends}
  \EndWhile
\end{algorithmic}
\end{algorithm}

The ChASE algorithm Alg.~\ref{alg:ChASE} is a more sophisticated
version of Rutishauser subspace iteration. The initial
input requires only the matrix $A$ but can also accept a set of vectors
$\subit V{0}$ which span the desired subspace. When specified as input, $\subit V{0}$ represents approximate solutions (e.g., when solving for a sequence of eigenproblems). When not specified, $\subit V{0}$ are initialized with normal random vectors ({\tt line~\ref{lst:line:init}}). The order $\ell$ of $\subit V{0}$
is such that the search space is larger (but not much larger) than the
desired number of eigenpairs $\nev$. 


%It is important to stress that while
%the full searched space is filtered by the Chebyshev polynomial, we
%are concerned only with the convergence properties of the subspace
%spanned by $k$ vectors $\subit{V_1}{i} \subset \subit{V}{i}$. 

%As mentioned at the beginning of Section~\ref{sec-2-1}, in opposition
%to Algorithm~\ref{alg:SubIt}, the desired part of the eigenspectrum is
%the lowest one. Notice that, thanks to the properties of the Chebyshev
%filter, solving for the highest part of the spectrum corresponds to
%just substitute $A$ with $-A$.  

In order to work efficiently the
Chebyshev filter requires estimates for $\lambda_1, \lambda_\ell$ and
$\lambda_n$. Only the latter has to be really bounded from above; the
first two numbers can be estimated without compromising the
effectiveness of the algorithm. All three values can be provided as
part of the input or can be computed by running few Lanczos iterations
({\tt line~\ref{lst:line:lanczos}}). Each time the body of the while
loop is repeated, the vectors $\subit V{i}$ are filtered ({\tt
  line~\ref{lst:line:cheby}}) by a Chebyshev polynomial with variable
degree $m_i$. The resulting vectors are orthonormalized using a
QR algorithm ({\tt line~\ref{lst:line:QR}}): this is the step which is the focus of this manuscript.  
The orthonormalization is followed by a
Rayleigh-Ritz projection ({\tt line~\ref{lst:line:rrstarts}}) at the
end of which eigenvector residuals are computed,
converged eigenpairs are deflated and locked ({\tt
  line~\ref{lst:line:dlstarts}}) while the non-converged vectors are
sent again to the filter to repeat the whole procedure ({\tt
  line~\ref{lst:line:dlends}}).
%\subsection{Locking, Deflation and Orthonormalization}
This last step ensures that eigenpairs that already satisfy the stopping criterion
are not further filtered which would unnecessarily increase the 
number of floating-point operations (FLOPs) performed by ChASE. 
%Therefore, 
%it is good practice to remove the converged eigenvectors and eigenvalues 
%from the arrays $V^{(i)}$ and $\Lambda^{(i)}$, and store them in the corresponding
%output arrays $Y$ and $\Lambda$ ({\tt line~\ref{lst:line:dlstarts}-\ref{lst:line:dlends}}). 
Only the remaining non-converged vectors are passed to the filtering step for
further refinement.

The orthonormalization phase is a key step of the algorithmic process. 
In the Hermitian case, ChASE not only orthonormalizes
the current search subspace but also ensures orthogonality against all previously
converged eigenvectors. This step enhances convergence by removing components 
of the space that correspond to already converged directions. Consequently, in 
ChASE, the QR factorization is always applied to a matrix of size $n\times \ell$
regardless of how many eigenpairs have already converged and been locked.
For numerical robustness, this QR factorization is conventionally performed 
using Householder QR implementation provided by LAPACK or ScaLAPACK. Numerical
stability in this step is crucial for the convergence behavior of ChASE, particularly 
during the initial iterations when the Chebyshev-filtered vectors can be 
highly correlated. While stability is an important ingredient, the ChASE library pays a price in terms of  performance due to the use of a less than ideal algorithm (Householder QR).
This is compounded by the observation that the QR factorization becomes increasingly
important as more eigenpairs converge and are locked during the iteration process.
When computing more than a small fraction of the spectrum, the orthonormalization
step can thus become a critical performance bottleneck \cite{Wu2022chase}.


\subsection{Replacing Householder QR with Communication-Avoiding CholeskyQR Variants}\label{sec:Replacing Householder QR with Communication-Avoiding CholeskyQR Variants}

On distributed memory architectures, and particularly on GPUs, the traditional Householder QR factorization often becomes a major performance bottleneck. Although Householder QR provides excellent numerical stability, its performance is primarily constrained by communication rather than computation. The algorithm requires frequent synchronization and numerous small data exchanges between processes, especially during panel factorizations. These synchronization points introduce latency that scales poorly with the number of processes and are even more detrimental on GPU-based systems, where the deep memory hierarchy and kernel-launch overheads amplify the cost of fine-grained communication. Consequently, as the number of locked eigenpairs increases in ChASE and the subspace to be orthonormalized grows, the cumulative cost of Householder QR can dominate the overall runtime \cite{Wu2022chase}.

To alleviate this bottleneck, we replaced the traditional Householder QR with communication avoiding (CA) variants of CholeskyQR~\cite{fukaya2014choleskyqr2,fukaya2020shifted}, which are more amenable to distributed and GPU-based architectures. In each iteration of ChASE, a QR factorization is performed on a rectangular matrix $\subit V{i} \in \mathbb{C}^{n \times k}$ with $n > k$. Among the available CAQR methods, we favor CholeskyQR over the tall-skinny QR (TSQR)~\cite{demmel2008communication}, even though both exhibit similar asymptotic communication costs. The main advantage of CholeskyQR lies in its simplicity: it relies on an additive reduction operation, while TSQR requires a QR factorization of a small matrix at every reduction step~\cite{fukaya2014choleskyqr2}. This difference simplifies the implementation and improves performance on hierarchical architectures, especially when efficient collective operations such as \texttt{MPI\_Allreduce} are available. Moreover, TSQR provides a clear performance advantage over ScaLAPACK’s Householder QR primarily when $n \gg k$~\cite{ballard2015reconstructing}, which is not typically the case in ChASE. For instance, in condensed-matter simulations, ChASE often targets about 10\% of the spectrum, leading to a moderate rather than an extreme $n/k$ ratios.

Although CholeskyQR delivers superior performance and scalability, it suffers from a loss of orthogonality when the condition number of $\subit V{i}$ increases. This loss can be mitigated by performing the procedure twice, yielding the CholeskyQR2 algorithm~\cite{fukaya2014choleskyqr2} (see Algorithm \ref{alg:cholqr}). However, CholeskyQR2 remains limited by the requirement that the Cholesky factorization of the Gram matrix $R = X^H X$ succeeds, which only holds when the condition number $\kappa_2(X)$ does not exceed $\mathcal{O}(\mathbf{u}^{-1/2})$, where $\mathbf{u}$ is the unit round-off and $\kappa_2(X) = \frac{\sigma_{\max}(X)}{\sigma_{\min}(X)}$.

\begin{algorithm}[htbp]
\caption{CholeskyQR for $X = QR$} \label{alg:cholqr}
\begin{algorithmic}[1]
  \Function{CholeskyQR}{$X$, {\sf cholDeg}}
    \For{$i = 1, \dots, ${\sf cholDeg}}
      \State $R \leftarrow X^H X$
      \State $[R, {\sf info}] \leftarrow \textsc{Cholesky}(R)$
      \If{${\sf info} \neq 0$}
        \State \Return {\sf Failure}
      \EndIf
      \State $X \leftarrow X R^{-1}$
    \EndFor
    \State \Return $X$
  \EndFunction
\end{algorithmic}
\end{algorithm}


To overcome this limitation, Fukaya et al.~\cite{fukaya2020shifted} proposed the shifted CholeskyQR2 (or $s$-CholeskyQR2) algorithm, which introduces a preconditioning step that shifts the Gram matrix $R$ by a small multiple of the identity. This shift effectively reduces the condition number of $X$, extending the applicability of CholeskyQR2 to matrices with $\kappa_2(X)$ up to $\mathcal{O}(\mathbf{u}^{-1})$.
Table~\ref{tab:cond_for_qr} summarizes the condition number ranges for each of the QR variants. This characterization enables the dynamic selection of the most appropriate QR factorization variant to balance performance and numerical accuracy.

\begin{table}[htbp]
\footnotesize
\renewcommand{\arraystretch}{1.}
\caption{Maximal acceptable condition numbers for different QR variants. $\mathbf{u}$ denotes the unit round-off.}
\label{tab:cond_for_qr}
\centering
\begin{tabular}{c|c|c|c|c}
\toprule
& CholeskyQR & CholeskyQR2 & $s$-CholeskyQR2 & HHQR \\
\midrule
$\kappa_2(X)$ & $\mathcal{O}(1)$ & $\mathcal{O}(\mathbf{u}^{-1/2})$ & $\mathcal{O}(\mathbf{u}^{-1})$ & Any \\
\bottomrule
\end{tabular}
\end{table}

Given its favorable performance characteristics, employing CholeskyQR and its variants within the ChASE library is highly attractive. At the same time, preserving numerical accuracy is essential and cannot be compromised. Achieving both performance and stability hinges on obtaining a reliable estimate of the condition number of the filtered vectors, $\kappa_2(V^{(i)})$. Directly computing this quantity is prohibitively expensive and would negate the performance gains offered by CholeskyQR. In this work, we demonstrate an alternative approach that provides a lightweight yet reliable estimate of the condition number without incurring significant overhead. We first present the resulting strategy and explain how it is used to adapt the QR step within the algorithm. The remainder of the paper is devoted to a rigorous derivation of these results and to numerical experiments that illustrate their impact on both the stability and performance of the ChASE library.


\subsubsection{Condition number estimates and dynamic CAQR}

For the simplest case when ChASE converges in one single step and does not optimize the polynomial degree (one degree for all vectors), we bound the condition number of the filtered vectors $\subit V{i}$ such that
\begin{equation}
\label{eq:simple-est}
\cd \left(\subit V{i} \right) \leq \eta \rto 1{m},
\end{equation}
where $\rto 1{m}$ is the convergence ratio of the smallest eigenpair (see Definition \ref{def:conv_ratio}). When polynomial degree optimization is enabled, the filtering degree is adapted for each single column of the array $\subit V{i}$. In this case, the exponent in $\rto 1{m}$ is substituted with the highest polynomial degree which typically corresponds to the last vector of the search space $m_\ell$. 
In the case when vectors are deflated and locked the expression for the bound gets further modified to
\begin{equation}
\label{eq:full-est}
\cd (\subit V{i}) \lesssim \eta \rto {k+1}{m_{k+1}} \rto 1{m_\ell-m_{k+1}}. 
\end{equation}
where $m_k$ is the Chebyshev polynomial degree corresponds to the Ritz pair of index $k$, and $\subit V{i} \gets \big[Y\; \subit V{i}\big]$ is the rectangular matrix to be factorized, including the converged eigenvectors $Y$
and newly filtered subspace .

These estimations rely solely on data already available within ChASE---namely, the spectral bounds $(c, e)$, the computed Ritz values $\Lambda$, the optimized filter degrees $m_i$, and the number of locked eigenpairs $k$. 
%\textsf{locked}. 
As already mentioned, these estimations are based on a detailed spectral analysis of the Chebyshev-filtered subspace and constitute the central contribution of this paper which is exposed in the following sections.
Based on the estimated condition number Eqs.~\eqref{eq:simple-est} and \eqref{eq:full-est}, we select the appropriate CholeskyQR variant following the heuristic in Algorithm~\ref{alg:caqr}. For double precision, if the estimated condition number exceeds $10^8$, the $s$-CholeskyQR2 variant is employed. For small condition numbers (below $20$), standard CholeskyQR suffices, while intermediate cases use CholeskyQR2. For robustness, a fallback mechanism (Algorithm~\ref{alg:caqr}, line~\ref{alg:line:scalapack_qr}) ensures that any failure in $s$-CholeskyQR2$ $ safely reverts to ScaLAPACK’s Householder QR routine.

\begin{algorithm}[htbp]
\caption{Dynamic CAQR selection in ChASE}\label{alg:caqr}
\begin{algorithmic}[1]
  \Function{Dynamic-CAQR}{$X$, \textsf{estCond}}
    \If{\textsf{estCond} $> 10^8$} \Comment{High condition number: use $s$-CholeskyQR2}
      \State $R \leftarrow X^H X$
      \State ${\sf norm} \leftarrow \|X\|_F$
      \State $s \leftarrow 11 (mn + n(n+1)) \, \mathbf{u} \, {\sf norm}$
      \State $[R, {\sf info}] \leftarrow \textsc{Cholesky}(R + sI)$
      \If{${\sf info} \neq 0$} \Comment{Cholesky failed, fallback to Householder QR}
        \State $X \leftarrow \textsc{HouseholderQR}(X)$ \label{alg:line:scalapack_qr}
      \Else
        \State $X \leftarrow X R^{-1}$
        \State $X \leftarrow \textsc{CholeskyQR}(X, \textsf{cholDeg}=2)$
      \EndIf
    \ElsIf{\textsf{estCond} $< 20$} \Comment{Well-conditioned: use single CholeskyQR}
      \State $X \leftarrow \textsc{CholeskyQR}(X, \textsf{cholDeg}=1)$
    \Else \Comment{Moderate conditioning: use CholeskyQR2}
      \State $X \leftarrow \textsc{CholeskyQR}(X, \textsf{cholDeg}=2)$
    \EndIf
    \State \Return $X$
  \EndFunction
\end{algorithmic}
\end{algorithm}

In section \ref{sec-2}, we describe the spectral properties of the filtered vectors $\subit V{i}$ both in the case of generic polynomial degree and when such a degree is optimized for each single filtered vector. Section \ref{sec-3} provides a numerical analysis of the bounds for $\cd (\subit V{i})$ is all possible cases. Section \ref{sec-4} is devoted to numerical tests of the condition number bounds, comparison of dynamic CAQR with the traditional QR decomposition and some scaling tests. In section \ref{sec-5} we summarize and conclude.  

%%% Local Variables: 
%%% mode: latex
%%% TeX-master: "main"
%%% End: